\section{Discussion}

The experiments separate forecast-origin correction from residual transition
learning. Against the preselected PGNN-style comparator, the open-loop lower
bootstrap endpoint clears the 5\% practical margin. In the simulated
receding-horizon study, the corresponding state-loss improvement is
\ClosedLoopImprovement\%, with all 15 system--condition means favoring
CFR-WM. The synchronous two-step outage increases CFR-WM error by 2.7\%,
versus 51.9\% for PGNN-style. These results consistently associate the
advantage under corrupted observations with the complementary context update.

The component evidence does not support the same interpretation for the
learned residual. Filtering followed by nominal propagation reaches 0.111
NRMSE-25, compared with 0.119 for CFR-WM, and removing the residual at
inference gives the lowest mean in every reported regime. Filter plus nominal
dynamics also records a slightly lower closed-loop state loss and lower action
effort than the complete model. For the present systems and fitted
checkpoints, the residual adds no demonstrated point-prediction or regulation
benefit.

The regime decomposition explains why a uniform gain is insufficient. The
corrupted and shifted regimes improve by \CorruptImprovement\%, whereas the
clean regime changes by \CleanImprovement\% and fails the prespecified
noninferiority criterion. Equation~\eqref{eq:optimal-gain} shows that equal
trust is optimal only for a particular relationship between prediction and
measurement errors. Channel-specific gains estimated from training or
validation innovations could adapt to unequal sensor reliability without
using test outcomes.

The measured-flight results further narrow the transfer claim. CFR-WM lowers
50-step position error in all four context conditions, while the matched
unfiltered residual model has lower full-state NRMSE in all eight condition--
horizon comparisons. CFR-WM's nominal 90\% coverage is 0.524--0.667 at 50
steps and is below the unfiltered model in all four conditions. A control
objective that emphasizes position can therefore rank models differently from
an equally weighted state metric, but the undercoverage precludes treating the
reported intervals as calibrated constraint margins.

The results suggest a modular design procedure for predictive control: first
validate the forecast-origin estimator, then retain a residual transition only
when it improves held-out rollout and closed-loop objectives. State weights
should follow the control specification, and joint constraints require a
covariance or calibrated prediction-set model that represents cross-state
dependence. The reported RTX 5070 latency remains a batched throughput
measurement; deadline compliance requires timing on the intended controller.
